
We consider the three steps of the modeling process.

\begin{description}
\item[Decision variables] The design of the cylinder depends on two
  variables, both expressed in centimeters:
  \begin{itemize}
  \item the radius of the basis: $r$,
   \item the height of the cylinder: $h$.
  \end{itemize}

  \begin{center}
\mytikz{tikz01.png}{
    \begin{tikzpicture}[thick]
  \node (a) [cylinder, shape border rotate=90, draw, minimum height=30mm, minimum width=15mm] {};
  \draw [<->] ([xshift=5pt]a.before bottom) -- ([xshift=5pt]a.after top) node [midway, right] {$h$};
\coordinate (center) at ($(a.before top)!0.5!(a.after top)$);
  \draw [<->] (a.before top) -- (center) node [midway] {$r$};
\end{tikzpicture}
}
  \end{center}

  
\item[Objective function] Since the thickness of the aluminium is the same
at any part of the can, the total surface of the cylinder has to be
  minimized. The objective function must provide the total surface as
  a function of the decision variables:
  \[
f:\R^2 \to \R: f(r,h).
  \]
  \begin{itemize}
\item Each basis is a circle of radius $r$, so its surface is $\pi
  r^2$.
 \item The side of the can is a rectangle of area $2\pi r h$.
  \end{itemize}
  Therefore, the objective function to minimize is
  \[
f(r,h) = 2 \pi r^2 + 2 \pi r h.
  \]
\item[Constraints] The volume of the can must be 0.33 liters, that is
  330 cm$^3$. The first constraint is therefore:
  \[
  \pi r^2 h = 330.
  \]
  We also need non negativity constraints:
  \begin{align*}
    r & \geq 0, \\
    h & \geq 0.
  \end{align*}
\end{description}

The optimization problem is therefore:

\[
\min_{r,h} 2 \pi r^2 + 2 \pi r h
\]
subject to
  \begin{align*}
  \pi r^2 h &= 330, \\
    r & \geq 0, \\
    h & \geq 0.
  \end{align*}

The optimal solution of this problem is $r = 3.746$ cm and $h = 7.491$ cm.


